% ============================================================
% sections/06_metodologia.tex — Metodologia
% ============================================================

\section{Metodologia}

Este trabalho adota uma abordagem de pesquisa aplicada com método de
desenvolvimento iterativo, combinando revisão bibliográfica sistemática,
design centrado no usuário e prototipação ágil \cite{mayer2019serious}.

\subsection{Classificação da Pesquisa}

Quanto à natureza: pesquisa aplicada, pois visa gerar um produto com
aplicação prática imediata. Quanto à abordagem: qualitativa e quantitativa
(mista). Quanto aos objetivos: exploratória e descritiva. Quanto aos
procedimentos: pesquisa bibliográfica e estudo de desenvolvimento.

% TODO: Fundamentar classificação metodológica com referência (ex: Gil, 2017)

\subsection{Etapas do Desenvolvimento}

O desenvolvimento do \textit{Alchemon} está organizado nas seguintes etapas:

\begin{enumerate}[1.]
    \item \textbf{Pesquisa bibliográfica}: revisão de literatura sobre jogos
          sérios, GBL e ensino de química;
    \item \textbf{Análise de requisitos}: levantamento das necessidades
          educacionais e técnicas;
    \item \textbf{Design do jogo}: elaboração do \textit{Game Design Document} (GDD);
    \item \textbf{Prototipação}: implementação do protótipo de baixa e alta fidelidade;
    \item \textbf{Avaliação} (TCC II): testes com usuários e análise de resultados.
\end{enumerate}

\subsection{Ferramentas e Tecnologias}

% TODO: Confirmar motor de jogo com orientador e atualizar esta subseção

\begin{table}[htbp]
    \centering
    \caption{Ferramentas Utilizadas no Desenvolvimento}
    \label{tab:ferramentas}
    \begin{tabular}{@{}lll@{}}
        \toprule
        \textbf{Ferramenta} & \textbf{Categoria} & \textbf{Versão} \\
        \midrule
        Unity               & Motor de Jogo      & 2022 LTS         \\
        Visual Studio Code  & Editor de Código   & 1.x              \\
        Blender             & Modelagem 3D       & 4.x              \\
        Git / GitHub        & Controle de Versão & ---              \\
        \bottomrule
    \end{tabular}
\end{table}

A Tabela~\ref{tab:ferramentas} lista as principais ferramentas adotadas no
projeto, selecionadas por serem gratuitas, amplamente utilizadas na indústria
e compatíveis com o ambiente de desenvolvimento do autor.

\subsection{Métricas de Avaliação}

A eficácia educacional do \textit{Alchemon} será mensurada por meio do
coeficiente de ganho normalizado de Hake:

\begin{equation}
    \langle g \rangle = \frac{\langle S_{pós} \rangle - \langle S_{pré} \rangle}
                             {100\% - \langle S_{pré} \rangle}
    \label{eq:hake}
\end{equation}

\noindent onde $\langle S_{pré} \rangle$ e $\langle S_{pós} \rangle$ representam
as médias dos escores percentuais em pré-teste e pós-teste, respectivamente.
Valores de $\langle g \rangle > 0{,}3$ são considerados ganhos satisfatórios
% NOTA: O uso de {,} em modo matemático requer \usepackage{icomma} ou siunitx
\cite{gee2003videogames}.

% TODO: Detalhar protocolo de coleta de dados e comitê de ética (se aplicável)
